\scanpage{635}
\mspage{23}

L'axiome de décomposition est satisfait pour
tout prédérivateur de la forme $\unsure{\mathbb{D}|\underline{D}}_{M,W}$, \deleted{\illegible{1}}
\deleted{\unsure{$W$}} l'axiome d'exactitude à gauche \inserted{Der 5 bis} est valable,
pourvu que les $(M(I),W(I))$ \unsure{s'inscrivent|s'insèrent} dans
une structure à cofibrations, ce qui est le cas.
De $=$ pour l'axiome Der 2 de localisation,
pourvu que de plus $W$ soit fortement
saturé, ce qui est également le cas. Enfin,
l'axiome de l'objet final dit que si
$I$ admet un objet final, alors

\[
 \operatorname{Ho}_W(e)\unsure{\;-\;|\longrightarrow}\operatorname{Ho}_W(I)
\]

est pl. fidèle, ou encore que le morphisme
d'adjonction

\[
 u_!u^*(\unsure{\xi|\zeta})\longrightarrow \xi
\]

est iso. Mais en vertu de \unsure{1.12}, cas $3^\circ$,
dans ce cas $u^*$ \deleted{\unsure{transforme}}\inserted{applique} $M(e)_c\simeq M_c$ dans
$M(I)_c$, donc le calcul du couple de
foncteurs adjoints $u_!^W,u^*_W$ peut se faire en
utilisant les \unsure{restrictions} de $u_!,u^*$ à $M(I)_c$ et
$\unsure{M_c}$ respectivement. Comme le morphisme
d'adjonction $u_!^c u_c^*\unsure{\longrightarrow|\mapsto}\mathrm{id}$ est iso, il en
est donc de $=$ de $u_!^W u^*_W$ qui s'en
déduit en localisant, \unsure{qed|cqfd}.

2.5. \unsure{Pour} achever de construire $\operatorname{Ho}_W$ comme
un dérivateur sur Cat \deleted{tout} entier, il faut
% archive sheet number: 634
